\textbf{COLUNI 2018 -- Questão 06}\par\medskip Leia o seguinte excerto que exemplifica a linguagem coloquial do livro de Carolina Maria de Jesus:

\par

“15 de julho de 1955 Aniversário da minha filha Vera Eunice. Eu pretendia comprar um par de sapatos para ela. Mas o custo dos generos alimenticios nos impede a realização dos nossos desejos. Atualmente somos escravos do custo de vida. Eu achei um par de sapatos no lixo, lavei e remendei para ela calçar. Eu não tinha um tostão para comprar pão. Então eu lavei 3 litros e troquei com o Arnaldo. Ele ficou com os litros e deu-me pão. Fui receber o dinheiro do papel. Recebi 65 cruzeiros. Comprei 20 de carne. 1 quilo de toucinho e um 1 de açucar e seis cruzeiros de queijo. E o dinheiro acabou-se.” (p.11)

\par

Na estruturação sintática da linguagem no excerto acima, há predominância de:\begin{enumerate}
\item períodos simples e orações subordinadas, criando uma narrativa densa com uma forte crítica social.
\item períodos simples e orações coordenadas, traduzindo com realismo a forma como a autora via o mundo.
\item períodos longos com orações coordenadas e subordinadas, construindo o ponto de vista da autora.
\item períodos longos permeados de recursos expressivos, construindo poeticamente o espaço social.\end{enumerate}
